% \begin{todolist}
%     \item why the membrane: we know from previous research that we can improve CNT stability by covering them: 
%     \item membrane composition: PVC based; this means that the electrolyte cannot penetrate theoretically; in practice, the device works very well, so the electrolyte still acts as dielectric. One of the potential issues that we have encountered is that ions get trapped in the membrane, potentially causing disruptions in the signal. Nevertheless, if we let enough time at the end of the collection of a tranfer curve (aka 5 seconds), the ions have plenty of time to get remmoved from the ion traps within the membrane and the measurements are not affected by this phenomenon
%     \item membrane deposition optimization: the first important step to optimize the deposition was to keep the vial with the membrane in ice; this helped to maintain it liquid for a longer time, allowing us to precisely deposit it (without the ice, the membrane dries in a few seconds, which is not enough time to keep it as even as possible). Having more time gives us the advantage of spreading it more evenly, thus mitigating the coffee ring effect (which is still there, as evidenced by the profilometry conducted) and guaranteeing that we are able to cover the whole area that we need to cover. wAs for the amount to deposit: we first tried to deposit 10ul in one go; this led to the membrane spreading too much and being too thin, giving us some mixed results (some devices were more covered than others, and there was high variability in the performance, approximately 50\% were stabilized and 50\% were not). This was not satisfactory and we tried again by depositing a first layer of 8ul, let it dry 30 minutes and cover it with another 7ul (total 15ul of membrane). This led to a more reliable performance, and we can say that the membrane effectively covers and protects the CNTs, as evidenced in the next section
%     \item membrane stability results (resistance, transfer curves) [paraphrase and discuss more in depth the next item point (it's from one of my own papers, so the results will be the same, but I need to change it slightly to not plagiarize myself)]
%     \begin{itemize}
%         \item 60 minutes
%         \item 4 hours
%         \item 12 hours
%         \item repeated
%     \end{itemize}
%     \item The performance evaluation of the devices before and after the membrane deposition was carried out by recording output and transfer characteristics in 1x PBS. The output charac- teristics showed distinct linear and saturation regions for the membrane free device, while for the encapsulated device the linear and saturation regimes were less distinguishable (Figure S5). One reason for this could be that the increased RDS, due to the presence of the membrane, limits the output current, thus preventing the device to reach saturation regime.
%     A thorough characterization in terms of stability and sta- bilization time took place through the collection of transfer characteristics. In Figure 2a it is possible to observe the results of a device without a membrane, displaying that after each of the 40 transfer cycles, the current decreases. To better observe the current trend in the devices, the ION was also extracted, i.e., the one measured at VGS=−0.8 V, and plotted over time: in Figure 2c, it is possible to see that the current drops drastically over the course of an hour, i.e.,−57.91μA after 5 minutes down to −22.50 μA after 60 minutes. This current decrease can be ascribed to the degradation of the SWCNTs when they come in contact with air and electrolyte. On the other hand, the devices with membrane encapsulation exhibit a different behaviour, showing that the current increases after each round of transfer (as seen in Figure 2b). Extracting the ION (Figure 2d), it is possible to notice an initial period in which the current drops immediately after the first measure- ment, and then rises again in a linear way. In this plot two phases can therefore be described: the first, previously defined by [39] as “stabilization phase”, involves the establishment of equilibrium in the electric field due to the rearrangement of ions in solution and within the membrane; the second, referred to as “constant slope phase”, is visible from 30 to 60 minutes, and is that phase in which the current changes linearly over time. In the example shown, there is an increase of 11 nA/min with a coefficient of determination of 0.99\%, as visible in Figure 2d. This phase is of particular importance for the calculation of the baseline that needs to be substracted from the measured IDS to obtain the corrected response, as in the example shown later of the EG-CNTFET as a signal transducer for a sensor. For a more detailed and in-depth description of the process, please refer to the paper by Petrelli et al. [39]. The linearity of the constant slope phase was evaluated with a fitting (linear regression analysis) and a device that showed R2 > 0.98 was deemed linear. The time necessary to reach this stage was used to quantify the differences between the devices and to evaluate any improvements compared to the literature. From the data collected it emerges that the average stabilization time of the encapsulated devices is 33.8 ± 12.9 minutes, N = 6, effectively halving the value obtained by Molazemhosseini et al. [33]. It is important to notice that while indeed the stabilization time is reduced and the stability improved, there is a price to pay in terms of current intensity: the IDS for the encapsulated devices is lower by one order of magnitude compared to the bare ones. This is due to the insulating properties of the membrane which, by increasing the resistance between the electrodes, limits the movement of charges and thus reduces the measured current. Although the device with the membrane exhibits a current that is an order of magnitude lower than the bare device, this reduced current level is still sufficiently high to ensure accurate measurements by our instruments. Therefore, we accept this reduction in current and focus on the additional benefits that the membrane offers.
%     Another important parameter to be evaluated for EG-FET- based devices is the VTH. Here the VTH is calculated using the extrapolation in the linear region method [40, 41], by fitting a linear function to the transfer curves at their point of maximum first derivative. The VTH is then found at the x-axis intercept of this linear fitting. In Figure 3a, it is reported the average VTH for devices without a membrane (N = 3) and devices with a membrane (N = 6) in red and blue, respectively. From this plot it is observed that this parameter keeps increasing in modulus for bare SWCNTs devices, whereas for the encap- sulated devices from 34.5 minutes the VTH shows a linear increase corroborating the results obtained in terms of IDS.
%     Figure 3b shows the average time change of the ION/IOFF of devices with and without the membrane. The bare devices show a decreases of the ION/IOFF starting from the initial value of 99.94 to a stabilized level of 61.18, representing a drop by 31\%, and indicating a deterioration in device performance. This decrease suggests that the ability of the device to switch between its on and off states is diminishing, leading to reduced sensitivity or reliability in its operation. In contrast, devices with membrane-encaspulated SWCNTs show a continuous linear increase over time; at the beginning of the measurement the ION/IOFF amounts to 9.77, then continues to increase to 24.51; the ×2.5 increase compared to the initial value indicates an improvement in performance, especially in terms of switching capabilities, that can be leveraged for measurements that need to be acquired over prolonged periods of time.
%     To further prove the sustained performance of the devices, measurements were carried out for 4 hours (Figure 4a) and 12 hours (Figure S6). Even in these cases, an initial stabilization phase is present, as previously observed, followed by a linear phase which is maintained for the entire duration of the mea- surement, i.e., an increase of 1 nA/min with a coefficient of determination of 99\%, allowing us to state that the membrane makes the devices usable for a prolonged time, without causing problems to the measurements.
%     An additional test that was carried out involved multiple acquisitions of 40 cycles of transfer characteristics on the same device. Figure 4b shows the example of a single device on which 5 consecutive measurements were carried out. Although there is a notable difference in the first curves measured, amounting to 9.5\%, it can be seen how the differences diminish over time, so much so that after 30 minutes, the variability between the 5 measurements is close to 2\%, a further confir- mation how the membrane encapsulation enhances the device stability. [I have to insert the figures, this will be done in a later stage of the thesis writing (aka, this afternoon)]
%     \item [These are the captions of the figures, I want you to incorporate them as best in the discussion of the results]:
%     figure 1: EG-CNTFET transfer characteristics collected over time, for devices without (a) and with (b) the membrane. Figure (a) shows that the current drops after each recorded cycle, while in (b) the opposite behavior is observed. (c) and (d) show how the normalized current ION/I0ON varies over time in the device without membrane and with membrane encapsulation, respectively. From these graphs, the difference in the behavior of the two devices is easier and more readily understood: in (c) it can be seen that the current drops very rapidly over the course of an hour, and the trend is set to reach 0. In (d) it can be seen that after an initial sudden drop, the current increases over time until it almost remains stable.
%     figure 2: The variation of two main EG-CNTFET parameters is evaluated: (a) shows the average variation of VTH for 3 devices without membrane and 6 devices with the membrane. As shown, there is a clear difference in the behaviour, namely the VTH of the devices without the membrane continues to decrease over time, indicating a worsening performance. The opposite is true for the encapsulated devices: the VTH keeps decreasing in modulus over time, reaching values lower than those of the bare devices already after 30 minutes (−0.581V for the bare devices and −0.576V for the encapsulated ones). The trend in (b) is slightly different, with the ION/IOFF of the devices without the encapsulation decreasing in a similar manner to the VTH; the ION/IOFF of the devices with the membrane shows an almost linear increase over the span of one hour, from 17.63 to to 24.51. 
%     figure 3: The variation of two main EG-CNTFET parameters is evaluated: (a) shows the average variation of VTH for 3 devices without membrane and 6 devices with the membrane. As shown, there is a clear difference in the behaviour, namely the VTH of the devices without the membrane continues to decrease over time, indicating a worsening performance. The opposite is true for the encapsulated devices: the VTH keeps decreasing in modulus over time, reaching values lower than those of the bare devices already after 30 minutes (−0.581V for the bare devices and −0.576V for the encapsulated ones). The trend in (b) is slightly different, with the ION/IOFF of the devices without the encapsulation decreasing in a similar manner to the VTH; the ION/IOFF of the devices with the membrane shows an almost linear increase over the span of one hour, from 17.63 to to 24.51.
% \end{todolist}

\gpt{
One of the primary challenges in employing CNT-based EG-FETs is their inherent instability when exposed to ambient conditions. Previous research has demonstrated that encapsulating CNTs can significantly enhance their stability, mitigating the degradation effects caused by oxygen and moisture. Based on these findings, we explored the use of a polymeric membrane as a protective layer to improve the long-term performance of our EG-CNTFETs.  
%
The encapsulating membrane used in this study is composed of a polyvinyl chloride (PVC)-based formulation. Theoretically, this composition prevents the penetration of the electrolyte; however, experimental results confirm that the device maintains its functionality, indicating that the electrolyte still effectively acts as a gate dielectric. Despite the improved stability conferred by the membrane, a potential issue arises from ion trapping within the polymeric matrix. This phenomenon can lead to transient disruptions in the signal; however, by introducing a \SI{5}{\second} delay at the end of each transfer curve measurement, we ensured sufficient time for trapped ions to dissipate, thereby preventing any significant impact on device performance.  
%
Optimization of the membrane deposition process was critical to achieving uniform coverage and minimizing device-to-device variability. A key refinement involved maintaining the membrane solution on ice during deposition, effectively prolonging its liquid state. This additional working time facilitated more precise deposition, mitigating the coffee-ring effect, as confirmed by profilometry measurements, and ensuring complete coverage of the CNTs. Initially, a single-step deposition of \SI{10}{\micro\liter} was tested, but this resulted in inconsistent thickness and variable device performance, with only approximately 50\% of devices exhibiting the desired stabilization. To address this, a two-step deposition strategy was employed: first, \SI{8}{\micro\liter} of membrane solution was deposited and allowed to dry for \SI{30}{\minute}, followed by an additional \SI{7}{\micro\liter} layer. This approach, totaling \SI{15}{\micro\liter} of membrane material, led to a significantly more reliable performance, effectively encapsulating the CNTs and enhancing their stability.  
%
The impact of membrane encapsulation on device stability was assessed through resistance measurements and repeated transfer curve acquisitions over extended periods. Transfer characteristics were recorded at intervals of \SI{60}{\minute}, \SI{4}{\hour}, and \SI{12}{\hour}, as well as over repeated measurement cycles. The performance evaluation was carried out by recording output and transfer characteristics in a \SI{1}{\times} PBS solution. The output curves revealed a clear distinction between membrane-free and encapsulated devices: while non-encapsulated EG-CNTFETs exhibited well-defined linear and saturation regimes, the encapsulated devices showed less distinguishable regions, likely due to an increase in $R_{\text{DS}}$ caused by the insulating nature of the membrane, which limited the output current and prevented full saturation.  
%
A thorough stability characterization was conducted by monitoring the evolution of transfer characteristics over time. Figure 1(a) shows that for non-encapsulated devices, the current decreases progressively with each transfer cycle, indicative of CNT degradation upon exposure to air and electrolyte. This trend is further corroborated by the extracted $I_{\text{ON}}$ values, plotted in Figure 1(c), which reveal a significant decline from \SI{-57.91}{\micro\ampere} after \SI{5}{\minute} to \SI{-22.50}{\micro\ampere} after \SI{60}{\minute}. In contrast, the encapsulated devices (Figure 1(b)) exhibited a markedly different behavior, with the current increasing after each cycle. The extracted $I_{\text{ON}}$ values, shown in Figure 1(d), highlight two distinct phases: an initial stabilization phase, attributed to ion rearrangement within the solution and membrane, followed by a linear increase in current over time. The latter, referred to as the "constant slope phase," is particularly significant as it defines the baseline correction needed for sensor applications. A linear fit of this phase, shown in Figure 1(d), yielded an increase rate of \SI{11}{\nano\ampere\per\minute}, with a coefficient of determination of \SI{0.99}{}. The average stabilization time for encapsulated devices was found to be \SI{33.8 \pm 12.9}{\minute} ($N=6$), effectively halving the time reported by Molazemhosseini et al.  
%
While membrane encapsulation substantially improved stability, it came at the cost of reduced current intensity due to the insulating nature of the polymeric layer. The drain current ($I_{\text{DS}}$) for encapsulated devices was lower by an order of magnitude compared to bare devices. However, this reduced current remained within a measurable range, ensuring reliable operation of the sensor.  
%
Threshold voltage ($V_{\text{TH}}$) variations were also analyzed to assess the impact of encapsulation on long-term device performance. Figure 2(a) shows the average $V_{\text{TH}}$ evolution for membrane-free and encapsulated devices. In non-encapsulated devices, $V_{\text{TH}}$ continuously increased in modulus over time, indicating worsening performance. Conversely, encapsulated devices exhibited a distinct stabilization behavior, where $V_{\text{TH}}$ decreased in modulus after approximately \SI{34.5}{\minute}, further confirming the membrane’s stabilizing effect. Similarly, the ratio of $I_{\text{ON}}/I_{\text{OFF}}$ (Figure 2(b)) displayed a declining trend for non-encapsulated devices, starting at \num{99.94} and stabilizing at \num{61.18}, representing a \SI{31}{\percent} reduction in performance. In contrast, encapsulated devices exhibited an almost linear increase in $I_{\text{ON}}/I_{\text{OFF}}$, from an initial value of \num{9.77} to \num{24.51} over the course of one hour, indicating improved switching performance.  
%
To further validate the long-term stability enhancement, measurements were extended to \SI{4}{\hour} (Figure 3(a)) and \SI{12}{\hour} (Figure S6). Even over these extended durations, the initial stabilization phase was followed by a sustained linear regime, with a stable increase of \SI{1}{\nano\ampere\per\minute} and a coefficient of determination of \SI{0.99}{}, confirming that the membrane effectively protects the CNTs and maintains device operability over prolonged periods.  
%
Additionally, the repeatability of measurements was assessed by conducting five consecutive sets of \num{40} transfer cycles on the same device (Figure 3(b)). While an initial deviation of \SI{9.5}{\percent} was observed, this difference diminished over time, with variability reducing to approximately \SI{2}{\percent} after \SI{30}{\minute}. This further corroborates that membrane encapsulation significantly enhances device stability and repeatability.  
%
Overall, the results presented in this section demonstrate that membrane encapsulation serves as an effective strategy to enhance EG-CNTFET stability, extending their operational lifetime while maintaining reliable performance. Although the reduced current levels introduce some trade-offs, the benefits of increased reproducibility, reduced drift, and improved $V_{\text{TH}}$ and $I_{\text{ON}}/I_{\text{OFF}}$ characteristics outweigh the drawbacks, making this approach highly suitable for sensor applications.
}
